\documentclass[11pt]{article}
\usepackage[utf8]{inputenc}
\usepackage[T1]{fontenc}
\usepackage[margin=1in]{geometry}
\usepackage{amsmath, amssymb, amsthm}
\usepackage{mathtools}
\usepackage[dvipsnames]{xcolor}
\usepackage{enumitem}
\usepackage{titlesec}

% Shortcuts
\newcommand{\R}{\mathbb{R}}
\newcommand{\N}{\mathbb{N}}
\newcommand{\eps}{\varepsilon}
\newcommand{\norm}[2]{\left\|#1\right\|_{#2}}

% Theorem environments
\newtheoremstyle{myplain}{8pt}{6pt}{\itshape}{}{\bfseries}{.}{6pt}{}
\theoremstyle{myplain}
\newtheorem{theorem}{Theorem}
\newtheorem{lemma}[theorem]{Lemma}
\newtheorem*{theorem*}{Theorem}

\theoremstyle{definition}
\newtheorem{example}{Example}
\newtheorem{definition}{Definition}

\theoremstyle{remark}
\newtheorem*{remark}{Remark}

% Technique highlight box
\newenvironment{technique}{%
  \par\smallskip
  \noindent\colorbox{Apricot}{\textbf{Technique}}\enspace\ignorespaces
}{\par\smallskip}

\titleformat{\section}{\large\bfseries}{}{0pt}{}
\titleformat{\subsection}{\normalsize\bfseries}{}{0pt}{}

\setlength{\parindent}{0pt}
\setlength{\parskip}{6pt}

\begin{document}

\begin{center}
  {\LARGE\bfseries Applying H\"{o}lder's Inequality}\\[4pt]
  {\large Math 104 $\cdot$ Spring 2026}
\end{center}

\bigskip

%======================================================================
\section{Statement}
%======================================================================

\begin{theorem*}[H\"{o}lder's Inequality]
Let $p, q > 1$ with $\dfrac{1}{p} + \dfrac{1}{q} = 1$ (conjugate exponents). If $f, g : [a,b] \to \R$ are Riemann integrable, then
\[
  \int_a^b |f(x)\,g(x)|\,dx \;\leq\; \norm{f}{p}\,\norm{g}{q}
  \;:=\; \left(\int_a^b |f|^p\right)^{1/p}\!\left(\int_a^b |g|^q\right)^{1/q}.
\]
The discrete version: for sequences $a_k, b_k$,
\[
  \sum_{k=1}^n |a_k b_k| \leq \left(\sum_{k=1}^n |a_k|^p\right)^{1/p}\!\left(\sum_{k=1}^n |b_k|^q\right)^{1/q}.
\]
\end{theorem*}

\begin{remark}
The case $p = q = 2$ is the \textbf{Cauchy--Schwarz inequality}. Young's inequality $uv \leq u^p/p + v^q/q$ underlies the proof: normalize $u = |f|/\norm{f}{p}$, $v = |g|/\norm{g}{q}$, apply Young pointwise, then integrate.
\end{remark}

%======================================================================
\section{Application 1: Bounding an Integral of a Product}
%======================================================================

\textbf{Setup.} You want to show $\int_a^b |h(x)|\,dx < \infty$ or bound it, where $h(x) = f(x)g(x)$ and neither factor alone is obviously integrable.

\begin{technique}
Split $h = f \cdot g$. Choose conjugate $p,q$ so that $\int |f|^p$ and $\int |g|^q$ are both finite. Then $\int|h| \leq \norm{f}{p}\norm{g}{q}$.
\end{technique}

\begin{example}
Show that $\displaystyle\int_0^1 \sqrt{x}\,e^{-x}\,dx$ is finite and bounded.

Take $f(x) = \sqrt{x} = x^{1/2}$ and $g(x) = e^{-x}$, with $p = 4$ and $q = 4/3$ (so $1/4 + 3/4 = 1$).
\[
  \int_0^1 \sqrt{x}\,e^{-x}\,dx \leq \left(\int_0^1 x^2\,dx\right)^{1/4}\!\left(\int_0^1 e^{-4x/3}\,dx\right)^{3/4}.
\]
Both factors are finite: $\int_0^1 x^2\,dx = 1/3$ and $e^{-4x/3}$ is bounded on $[0,1]$.
\end{example}

\begin{example}
Show that $\displaystyle\int_0^1 \frac{f(x)}{\sqrt{x}}\,dx$ is finite whenever $f \in L^2[0,1]$ (i.e., $\int_0^1 f^2\,dx < \infty$).

Apply Cauchy--Schwarz ($p = q = 2$) with $f(x)$ and $g(x) = 1/\sqrt{x}$:
\[
  \int_0^1 \frac{|f(x)|}{\sqrt{x}}\,dx \leq \left(\int_0^1 f(x)^2\,dx\right)^{1/2}\!\left(\int_0^1 \frac{1}{x}\,dx\right)^{1/2}.
\]
Wait --- $\int_0^1 x^{-1}\,dx$ diverges, so $p = q = 2$ fails here. Instead take $p = 4$, $q = 4/3$:
\[
  \int_0^1 \frac{|f(x)|}{\sqrt{x}}\,dx \leq \left(\int_0^1 |f|^4\right)^{1/4}\!\left(\int_0^1 x^{-2/3}\,dx\right)^{3/4}.
\]
Now $\int_0^1 x^{-2/3}\,dx = 3$ (since $-2/3 > -1$), so this works whenever $f \in L^4[0,1]$.

\textbf{Key lesson:} The choice of $p$ and $q$ matters. If one choice makes a factor diverge, try different exponents.
\end{example}

%======================================================================
\section{Application 2: The Trivial Split ($g \equiv 1$)}
%======================================================================

\textbf{Setup.} You want to bound $\norm{h}{1} = \int|h|$ using only information about $\norm{h}{p}$ for some $p > 1$.

\begin{technique}
Write $h = h \cdot 1$ and apply H\"{o}lder with $f = h$, $g \equiv 1$:
\[
  \int_a^b |h|\,dx \leq \left(\int_a^b |h|^p\right)^{1/p}(b-a)^{1/q}.
\]
This converts an $L^p$ bound into an $L^1$ bound, at the cost of a factor of $(b-a)^{1/q}$.
\end{technique}

\begin{example}
Let $f$ be continuous on $[0,1]$ with $\left(\int_0^1 f^2\right)^{1/2} \leq 5$. Bound $\int_0^1 |f|$.

Apply Cauchy--Schwarz with $g \equiv 1$:
\[
  \int_0^1 |f|\,dx \leq \left(\int_0^1 f^2\right)^{1/2}\!\left(\int_0^1 1^2\right)^{1/2} = 5 \cdot 1 = 5.
\]
\end{example}

\begin{example}[Comparing $L^p$ norms on finite intervals]
For $1 \leq r < s$ and $f \in L^s[a,b]$, show $\norm{f}{r} \leq (b-a)^{1/r - 1/s}\norm{f}{s}$.

Apply H\"{o}lder with $p = s/r > 1$ and $q = s/(s-r)$, writing $|f|^r = |f|^r \cdot 1$:
\[
  \int_a^b |f|^r \leq \left(\int_a^b |f|^s\right)^{r/s}(b-a)^{1-r/s}.
\]
Taking $r$-th roots: $\norm{f}{r} \leq \norm{f}{s}\cdot(b-a)^{1/r - 1/s}$.

\textbf{Moral:} On a finite interval, higher $L^p$ norms dominate lower ones (up to a constant depending on interval length).
\end{example}

%======================================================================
\section{Application 3: Cauchy--Schwarz ($p = q = 2$)}
%======================================================================

\textbf{Setup.} The most common special case. Use when both factors are naturally square-integrable or when symmetry suggests equal exponents.

\begin{technique}
$p = q = 2$: $\displaystyle\int |fg| \leq \left(\int f^2\right)^{1/2}\!\left(\int g^2\right)^{1/2}$.

In $\R^n$ (discrete): $\left|\sum a_k b_k\right| \leq \left(\sum a_k^2\right)^{1/2}\!\left(\sum b_k^2\right)^{1/2}$.
\end{technique}

\begin{example}
Prove that $\left(\int_a^b f(x)\,dx\right)^2 \leq (b-a)\int_a^b f(x)^2\,dx$.

This is the trivial split with Cauchy--Schwarz:
\[
  \left(\int_a^b f\right)^2 = \left(\int_a^b f \cdot 1\right)^2 \leq \left(\int_a^b f^2\right)\!\left(\int_a^b 1\right) = (b-a)\int_a^b f^2.
\]
\end{example}

\begin{example}
Prove the Cauchy--Schwarz inequality for inner products in $\R^n$: $|\langle u, v\rangle| \leq \|u\|\|v\|$.

Apply the discrete H\"{o}lder with $p=q=2$ to sequences $a_k = u_k$, $b_k = v_k$:
\[
  \left|\sum_{k=1}^n u_k v_k\right| \leq \sum_{k=1}^n |u_k||v_k| \leq \left(\sum u_k^2\right)^{1/2}\!\left(\sum v_k^2\right)^{1/2} = \|u\|\|v\|.
\]
\end{example}

%======================================================================
\section{Application 4: Bounding Improper Integrals (Comparison via H\"{o}lder)}
%======================================================================

\textbf{Setup.} You have an improper integral $\int_1^\infty f(x)/x^s\,dx$ and need to show it converges. If $f$ is bounded, comparison suffices; but if $f$ grows, H\"{o}lder can isolate the growth into a factor you can control.

\begin{technique}
Write $f(x)/x^s = f(x) \cdot x^{-s}$. Choose $p,q$ so that $\int |f|^p$ converges and $\int x^{-sq}$ converges (requires $sq > 1$, i.e., $s > 1/q = 1 - 1/p$).
\end{technique}

\begin{example}
Show $\displaystyle\int_1^\infty \frac{x-[x]}{x^{s+1}}\,dx$ converges for all $s > 0$.

Since $0 \leq x - [x] < 1$, we have $\left|\dfrac{x-[x]}{x^{s+1}}\right| \leq x^{-(s+1)}$. For $s > 0$, $s + 1 > 1$, so by comparison:
\[
  \int_1^\infty \frac{x-[x]}{x^{s+1}}\,dx \leq \int_1^\infty x^{-(s+1)}\,dx = \frac{1}{s} < \infty.
\]
(Here comparison suffices since the numerator is bounded. H\"{o}lder would be needed if $f$ grew.)
\end{example}

\begin{example}
Show $\displaystyle\int_1^\infty \frac{\ln x}{x^s}\,dx$ converges for $s > 1$.

Apply H\"{o}lder with $p = 2$, $q = 2$, writing $\ln x / x^s = (\ln x / x^{s/2}) \cdot (1/x^{s/2})$:
\[
  \int_1^\infty \frac{\ln x}{x^s}\,dx \leq \left(\int_1^\infty \frac{(\ln x)^2}{x^s}\,dx\right)^{1/2}\!\left(\int_1^\infty \frac{1}{x^s}\,dx\right)^{1/2}.
\]
Both factors converge for $s > 1$ (the second is $(s-1)^{-1/2}$; the first by integration by parts).
\end{example}

%======================================================================
\section{Application 5: Discrete H\"{o}lder for Finite Sums}
%======================================================================

\textbf{Setup.} You need to bound $\sum a_k b_k$ or $\sum |a_k|$ over a finite index set.

\begin{technique}
Apply the discrete version. The key choice: if you know $\ell^p$ information about one sequence (e.g., it comes from a norm), match it with the conjugate exponent on the other.
\end{technique}

\begin{example}
Prove that $\displaystyle\sum_{k=1}^n k \leq n^{3/2}\left(\sum_{k=1}^n k^2\right)^{1/2} / n^{1/2}$ --- actually, directly: $\sum_{k=1}^n k \leq \sqrt{n}\left(\sum_{k=1}^n k^2\right)^{1/2}$.

Apply Cauchy--Schwarz with $a_k = 1$, $b_k = k$:
\[
  \sum_{k=1}^n k = \sum_{k=1}^n 1 \cdot k \leq \left(\sum_{k=1}^n 1^2\right)^{1/2}\!\left(\sum_{k=1}^n k^2\right)^{1/2} = \sqrt{n}\left(\frac{n(n+1)(2n+1)}{6}\right)^{1/2}.
\]
\end{example}

\begin{example}
Show that for $x_1, \ldots, x_n \geq 0$: $\left(\dfrac{x_1 + \cdots + x_n}{n}\right)^p \leq \dfrac{x_1^p + \cdots + x_n^p}{n}$ for $p \geq 1$ (Jensen / power mean inequality).

By H\"{o}lder with $a_k = x_k$ and $b_k = 1/n$, exponents $p$ and $q = p/(p-1)$:
\[
  \frac{\sum x_k}{n} = \sum x_k \cdot \frac{1}{n} \leq \left(\sum x_k^p\right)^{1/p}\!\left(\sum \frac{1}{n^q}\right)^{1/q} = \left(\sum x_k^p\right)^{1/p} \cdot n^{1/q - 1}.
\]
Since $1/q - 1 = -1/p$: $\left(\sum x_k/n\right)^p \leq \sum x_k^p / n$.
\end{example}

%======================================================================
\section{Application 6: Moment Generating Functions (Probability)}
%======================================================================

\textbf{Setup.} The \textbf{moment generating function} (MGF) of a random variable $X$ is
\[
  M_X(t) = E[e^{tX}] = \int e^{tx}\,dP,
\]
where $dP$ is the probability measure. This is just a Lebesgue integral against a probability measure, so Hölder applies directly with $d\mu = dP$.

\subsection{Log-Convexity of the MGF}

\begin{theorem}
$M_X(t)$ is \textbf{log-convex}: $\log M_X$ is a convex function of $t$.
\end{theorem}

\begin{proof}
Let $\alpha \in (0,1)$, and set $p = 1/\alpha$, $q = 1/(1-\alpha)$, so $1/p + 1/q = 1$. Then:
\[
  M_X(\alpha s + (1-\alpha)t)
  = E\!\left[e^{\alpha s X} \cdot e^{(1-\alpha)tX}\right]
  \leq \bigl(E[e^{sX}]\bigr)^\alpha \bigl(E[e^{tX}]\bigr)^{1-\alpha}
  = M_X(s)^\alpha\, M_X(t)^{1-\alpha},
\]
where the inequality is Hölder with exponents $p = 1/\alpha$ and $q = 1/(1-\alpha)$.
Taking $\log$ of both sides gives
\[
  \log M_X(\alpha s + (1-\alpha)t) \leq \alpha \log M_X(s) + (1-\alpha)\log M_X(t),
\]
which is exactly convexity of $\log M_X$.
\end{proof}

\begin{technique}
To prove log-convexity of any quantity of the form $F(t) = E[f(t, X)]$ where $f(t,x) = e^{tx}$ or similar: write $F(\alpha s + (1-\alpha)t) = E[f(s,X)^\alpha f(t,X)^{1-\alpha}]$ and apply Hölder with $p = 1/\alpha$.
\end{technique}

\begin{remark}
The \textbf{cumulant generating function} (CGF) is $K(t) = \log M_X(t)$. Its first derivative at 0 is $E[X]$; its second derivative at 0 is $\mathrm{Var}(X)$. Convexity of $K$ means $K''(t) \geq 0$ everywhere, i.e., the ``running variance'' $\mathrm{Var}(e^{tX})$ is always nonneg --- a tautology, but Hölder is what proves it globally, not just at $t=0$.
\end{remark}

\subsection{Chernoff Bounds}

Log-convexity is the engine behind the sharpest tail bounds in probability.

\begin{example}[Chernoff bound]
For any $a \in \R$ and $t > 0$:
\[
  P(X \geq a) = P(e^{tX} \geq e^{ta}) \leq \frac{E[e^{tX}]}{e^{ta}} = e^{-ta} M_X(t).
\]
Optimizing over $t > 0$ gives the \textbf{Chernoff bound}:
\[
  P(X \geq a) \leq \inf_{t > 0}\, e^{-ta} M_X(t) = e^{-\sup_{t>0}(ta - \log M_X(t))}.
\]
The exponent is the \textbf{Legendre transform} of $K(t) = \log M_X(t)$. Convexity of $K$ (proved via Hölder above) guarantees the supremum is attained at a unique $t^*$ where $K'(t^*) = a$, making the bound explicit and tight.
\end{example}

\begin{example}[Large deviations / Cramér's theorem]
Let $S_n = (X_1 + \cdots + X_n)/n$ be the sample mean. The \textbf{rate function} governing how fast $P(S_n \geq a)$ decays is
\[
  I(a) = \sup_{t \in \R}\, (ta - K(t)),
\]
the Legendre--Fenchel transform of $K$. Since $K$ is convex (Hölder), $I$ is a proper convex function and equals zero only at $a = E[X]$. Cramér's theorem says $P(S_n \geq a) \approx e^{-nI(a)}$ --- so Hölder, via log-convexity, is literally the foundation of the exponential decay rate in the law of large numbers.
\end{example}

\subsection{The Abstract Picture: Convex Duality}

Both Hölder's inequality and MGFs are special cases of the same underlying structure in convex analysis.

Young's inequality --- the engine of Hölder's proof --- states:
\[
  uv \leq \frac{u^p}{p} + \frac{v^q}{q} \qquad \left(\tfrac{1}{p}+\tfrac{1}{q}=1\right).
\]
This is the \textbf{Fenchel--Young inequality} $\langle x, y\rangle \leq f(x) + f^*(y)$ applied to $f(u) = u^p/p$, whose convex conjugate is $f^*(v) = v^q/q$.

The Chernoff bound's optimization $\sup_t(ta - K(t))$ is exactly the convex conjugate $K^*(a)$ of the CGF. So the full chain is:

\begin{center}
\renewcommand{\arraystretch}{1.4}
\begin{tabular}{ccc}
  Young's inequality & $\longrightarrow$ & Hölder's inequality \\
  (Fenchel--Young, $f = u^p/p$) & & (bound $E[|fg|]$) \\[4pt]
  Hölder applied to $E[e^{tX}]$ & $\longrightarrow$ & Log-convexity of $M_X(t)$ \\[4pt]
  Legendre transform of $\log M_X$ & $\longrightarrow$ & Chernoff / Cramér rate function \\
\end{tabular}
\end{center}

\begin{remark}
At the deepest level, Hölder's inequality, MGFs, and large deviations are all facets of \textbf{convex duality}: the pairing between a convex function and its conjugate. Young's inequality is where it starts; the Cramér rate function is where it ends up.
\end{remark}

%======================================================================
\section{Decision Guide: Choosing $p$ and $q$}
%======================================================================

When you see $\int |fg|$ or $\sum |a_kb_k|$, ask:

\begin{enumerate}[leftmargin=*, itemsep=4pt]
  \item \textbf{Is the product symmetric?} Use $p = q = 2$ (Cauchy--Schwarz). Fastest and most natural.

  \item \textbf{Is one factor bounded?} Use comparison directly (H\"{o}lder is overkill). E.g., $|f| \leq M$ gives $\int |fg| \leq M\int|g|$.

  \item \textbf{Does one factor have a singularity like $x^{-\alpha}$?} Choose $q$ so that $q\alpha < 1$ (keeping the singular factor integrable). Then set $p = q/(q-1)$.

  \item \textbf{Do you only know an $L^p$ norm?} Use the trivial split with $g \equiv 1$ to get an $L^1$ bound. Result: $\int|h| \leq \norm{h}{p}(b-a)^{1/q}$.

  \item \textbf{Does a weight function appear?} Write $\int w(x)|f(x)|\,dx = \int |f| \cdot |w|$ and put the weight $w$ in the $L^q$ factor (often easier to compute) and $f$ in the $L^p$ factor (often the unknown).
\end{enumerate}

\begin{remark}
If your choice of $p,q$ makes one factor infinite, try different exponents or split the integral into $\int_0^1 + \int_1^\infty$ and use different exponents on each piece.
\end{remark}

\end{document}
